\section{Complexity and Numerical Behaviour}\label{sec:results}

\subsection{Complexity}\label{ssec:complexity}

Table~\ref{tab:complexity} collects the asymptotic cost of solving~\eqref{eq:bqp} by each
method, for $A = M^\top M$ with $M \in \mathbb{R}^{m\times n}$, condition number
$\kappa = \kappa(A_\alpha)$ after the regularising split, tolerance $\varepsilon$, free-set
size $k = |\mathcal{F}| \leq n$, and $s$ active-set outer steps. ``Extra memory'' is storage
beyond the operator; the matrix-free rows never form $A$ and need only $O(n)$ working memory.

\begin{table}[ht]
\centering
\footnotesize
\begin{tabular}{lcccc}
\toprule
Method & Per-step cost & Iterations & Extra memory & Forms $A$ \\
\midrule
Interior point (dense QP)          & $O(n^3)$        & $O(\log 1/\varepsilon)$          & $O(n^2)$ & Yes \\
Dense Cholesky $+$ active set       & $O(k^2 m)$      & $s$                             & $O(k^2)$ & Yes \\
CG matrix-free $+$ active set        & $O(km)$         & $s \cdot O(\sqrt\kappa)$        & $O(n)$   & No  \\
Factor / Woodbury $+$ active set     & $O(kr^2 + r^3)$ & $s$                             & $O(nr)$  & No  \\
Projected gradient (matrix-free)     & $O(nm)$         & $O(\kappa)$                     & $O(n)$   & No  \\
\bottomrule
\end{tabular}
\caption{Asymptotic cost of solving the non-negative quadratic~\eqref{eq:bqp} with
$A = M^\top M$, $M \in \mathbb{R}^{m \times n}$. Here $k = |\mathcal{F}| \leq n$ is the
free-set size, $s$ the number of active-set outer steps, $\kappa = \kappa(A_\alpha)$ the
condition number after the split~\eqref{eq:regsplit}, and $r$ the factor rank of a
diagonal-plus-low-rank operator (Section~\ref{sec:operators}). The last four rows share the
active-set loop and differ only in the backend~(Section~\ref{sec:operators}): matrix-free CG
and the Woodbury factor solve both avoid assembling $A$, the former iteratively at
$O(\sqrt\kappa)$ inner steps, the latter by an exact rank-$r$ inverse. The active-set rows differ only in the
inner solve: matrix-free CG applies $M_{\mathcal{F}}^\top(M_{\mathcal{F}}\,v)$ at $O(km)$ and
never assembles the Gram matrix, whereas dense Cholesky forms and factorises
$A_{\mathcal{F}}$ at $O(k^2 m)$; the matrix-free advantage grows with $k$. The CG
``Iterations'' column counts $s \cdot O(\sqrt\kappa)$ inner iterations in total; because the
loop starts at $\mathcal{F} = \{1,\dots,n\}$, the first solve at $k = n$ dominates and is
governed by the full $\kappa$ bounded in Proposition~\ref{prop:regcond}. Projected gradient
avoids the linear solve entirely but pays $O(\kappa)$, a factor $\sqrt\kappa$ more inner
iterations than CG at the same per-step cost.}
\label{tab:complexity}
\end{table}

Two comparisons are structural. Within the active-set family, matrix-free CG and dense
Cholesky share the outer loop and differ only in whether $A_{\mathcal{F}}$ is assembled; the
$O(km)$-versus-$O(k^2 m)$ gap means CG pulls ahead as the free set grows, while dense Cholesky
can win when $k$ is small or $\kappa$ is large enough that CG needs many inner iterations
(the crossover is near $\kappa \sim k^2$). Between the two matrix-free methods, CG's
$O(\sqrt\kappa)$ dominates projected gradient's $O(\kappa)$ whenever $\kappa$ is large; only
when the regularising split compresses $\kappa$ to $O(1)$ do the inner counts become
comparable, at which point projected gradient's loop-free simplicity is attractive.

\subsection{Behaviour on synthetic SPD problems}\label{ssec:synthetic}

The propositions pin down how the solver behaves on a controlled family of test problems.
Take $A = Q\Lambda Q^\top$ with $Q$ a random orthogonal matrix and a prescribed spectrum
$\Lambda = \mathrm{diag}(\lambda_1,\dots,\lambda_n)$ of target condition number $\kappa$
(equivalently $A = M^\top M$ for a random $M$ with the desired singular values), and plant a
non-negative solution $x^\star \geq 0$ with a chosen support size $|\{i : x^\star_i > 0\}|$ by
setting $b = A x^\star - s^\star$ for a complementary $s^\star \geq 0$. Three qualitative
predictions follow directly and are what the companion benchmarks confirm on their structured
operators:

\begin{itemize}
  \item \emph{Inner count tracks $\sqrt\kappa$.} By Proposition~\ref{prop:cg} the CG
  iterations to a fixed tolerance grow like $\sqrt\kappa$, so the regularising
  split~\eqref{eq:regsplit} -- which lowers $\kappa$ at rate $O((1-\alpha)/\alpha)$
  (Proposition~\ref{prop:regcond}) -- reduces the count steeply as $\alpha$ increases from
  zero. The effect is largest in the rank-deficient regime $m < n$, where the unsplit system
  is singular and only $\alpha > 0$ makes CG well posed.
  \item \emph{Outer count is small.} Block principal pivoting adjusts the free set by a few
  indices per step and, being non-cyclic through the Bland fallback
  (Theorem~\ref{thm:termination}), reaches the correct support in a number of outer steps that
  scales with how far the initial full set is from the optimal support, not with $n$; the
  patience budget is essentially never exhausted on generic data.
  \item \emph{CG beats projected gradient by $\sqrt\kappa$.} At matched per-step cost the two
  matrix-free methods separate by exactly the $O(\sqrt\kappa)$-versus-$O(\kappa)$ factor of
  Propositions~\ref{prop:cg} and~\ref{prop:proj}, the gap widening with $\kappa$ and closing
  under heavy regularisation.
\end{itemize}

Quantitative wall-clock studies on structured operators arising in applications -- including
scaling with problem size and warm-started parametric sweeps -- are reported in the companion
papers~\cite{schmelzer2026matrixfree,schmelzer2026rmt}, whose numerical sections instantiate
Algorithm~\ref{alg:elim} with the matrix-free CG inner solver analysed here.
